\section{Introduction}\label{sec:introduction51}
Let $\calA$ be a finite alphabet of orthogonal maps of $\R^D$.
For a word $w=a_1\cdots a_N$ write $U_w=U_{a_N}\cdots U_{a_1}$.
A seed $x$, a command word $w$, and a terminal coordinate query $j$
define the binary experiment
\begin{equation}\label{eq:target}
 \Pp(B=b\mid x,w,j)=\frac{1+b\rho e_j^TU_wx}{2},
 \qquad b\in\{-1,1\},\quad 0<\rho<1.
\end{equation}
A stochastic realization retains one finite register between commands.
Its transition row can depend on the current command and external
epoch, but not on a future command or on an unrecorded past. The
quantity $W_{N,\varepsilon}$ is the least peak number of available
register labels that approximates every prescribed binary law within
$\varepsilon$ in total variation. A different realization, including
its real transition rows, may be chosen for each horizon. Thus this
is a nonuniform clocked resource, not an autonomous computational
complexity measure.

The distinction between dimension of response and dimension of a
positive realization is classical. Invariant polyhedral cones describe
positive realizability for stationary systems \cite{BF04}; nonnegative
Hankel factorization alone does not impose a common positive
intertwining map \cite{BF03}. Here the issue is compatibility across
all commands and all cuts, even when the realization may change its
rows at every epoch. The signed linear response dimension is small,
but a state distribution must remain nonnegative, every transition
row must sum to one, and the same row must answer every future word.

For the structural results we prescribe the signed spanning seeds
$X=\{\pm e_1,\ldots,\pm e_D\}$, include all coordinate queries, and
assume $I\in\calA$. The normalized residual matrix then has ordinary
rank $D+1$. We call a cut \emph{rank-tight} when its available width
is exactly $D+1$. Ordinary rank is only a lower bound for positive
width; it does not assert a simultaneous stochastic factorization.

\subsection{Compatible lifts and bounded-width rigidity}
At arbitrary width, the observable state object is not generally the
convex hull of individual hidden responses. Let $R_t$ be the span of
actually reachable distributions, and let $Z_t$ record the actual
identity-suffix means. The relevant object is
\[
       P_t=R_t\cap\Delta_{K_t-1},\qquad S_t=(P_tZ_t)^T.
\]
Theorem~\ref{thm:section51} proves the command identities on $R_t$,
the positive inclusions between the sections, and the bound
$f_0(S_t)\le\binom{K_t}{\dim R_t-1}$. At one surplus label,
$K_t=D+2$, the two possible reachability ranks lead respectively to a
polytope with at most $D+2$ vertices or one with at most $D+2$ facets.
This distinction is lost by replacing the section with the whole
simplex.

Theorem~\ref{thm:local51} gives an exact finite characterization of
arbitrary prescribed profiles by orthogonal projectors, bounded mean
matrices and stochastic transitions. Its equations have degree at
most three and a description polynomial in the explicitly listed
epochs and register sizes, without listing all command words. This
is a real-algebraic existence and synthesis statement; it is not a
polynomial-time algorithm. When the sections and observable maps
are supplied, the remaining stochastic extension question has the
ordinary linear alternative of Proposition~\ref{prop:dual51}.

A compactness argument now applies at every bounded width, not only
at ordinary rank. Theorem~\ref{thm:occupation51} bounds the number
of cuts below any fixed width by a volume-expansion constant whenever
the command group is infinite. Arbitrarily larger intermediate
registers are permitted. Consequently:
\begin{quote}
For signed spanning seeds and all coordinate queries, the exact
clocked widths are bounded over all horizons if and only if the
orthogonal command group is finite. In that case a single finite
permutation machine works at every horizon.
\end{quote}
This is Theorem~\ref{thm:finitegroup51}. A uniform clocked bound $K$
gives a stationary realization with at most
$\binom K{\lfloor K/2\rfloor}$ labels; equality of these two resource
bounds is not asserted. A rational planar rotation of infinite order
together with a reflection provides a fixed-dimensional,
fixed-alphabet, noncommuting family with $W_{N,0}\to\infty$
(Corollary~\ref{cor:rational51}).

The one-surplus distinction has a further consequence not requiring
an infinite group. If $D\ge3$ and $I,-I\in\calA$, a full-dimensional
centrally symmetric polytope cannot have only $D+2$ vertices or
facets. Compactness makes this obstruction quantitative as a strictly
positive variational volume gap. Theorem~\ref{thm:surplus-budget51}
bounds the occupation of width at most $D+2$ and proves that in
dimension three a signed-permutation alphabet containing $I,-I$
has $W_{N,0}=6$ for all sufficiently large $N$. Its gap is
variational; no sharp numerical first horizon is claimed.

\subsection{Stable rank-tight occupation}
At a rank-tight cut the section is a simplex, and the classical
simplex difference-body identity yields, for $m$ selected command
cuts,
\begin{equation}\label{eq:intro-budget50}
 c_D^{m-1}\le D!\rho^{-D},\qquad
 c_D=2^{-D}\binom{2D}{D}>1,
\end{equation}
when $I,-I\in\calA$. The selected cuts may include $0$ and $N$;
intervening widths are unrestricted. This exact statement is retained
as Theorem~\ref{thm:occupation50}.

To pass to positive error, one must recover statewise information
without dividing by an uncontrolled hidden probability. The signed
seed distributions supply a uniformly conditioned affine anchor
matrix. Lemma~\ref{lem:anchor51} gives an explicit inverse bound and
recovers every hidden suffix row, with no sum of errors over epochs.
Writing $\delta=2\varepsilon$, $b=\rho-D\delta$, and
\[
 \Lambda=1+
 \frac{2D(D+1)(1+\sqrt D)\delta}
      {(\rho-2D\delta)(\rho-D\delta)},\qquad
                  0\le\delta<\frac{\rho}{2D},
\]
Theorem~\ref{thm:robust51} proves
\[
                 (c_D/\Lambda^D)^{m-1}\le D!b^{-D}.
\]
For planar signed-permutation alphabets containing $I,-I$, this gives
an error interval independent of the horizon:
\[
 0\le\varepsilon\le\rho^2/2000,\qquad
 (512/363)^N>2(500/499)^2\rho^{-2}
       \quad\Longrightarrow\quad W_{N,\varepsilon}=4.
\]
In particular $\rho=1/10$, $N\ge16$ and
$\varepsilon\le1/200000$ suffice. These constants are explicit
sufficient bounds. The earlier zero-error bound $N\ge14$ for this
signal is also retained; neither number is asserted to be the first
optimal horizon. The arbitrary-width finite-group theorem concerns
exact realization, whereas the displayed stability theorem concerns
rank-tight cuts. The hypotheses are not interchanged.

\subsection{Finite certificates and the rest of the theory}
The structural results complement, rather than replace, the exact
two-state sign--magnitude alternative and its algebraic synthesis.
The rational nine-dimensional, three-epoch orthogonal example of
Theorem~\ref{thm:orthogonal50} has a cubic optimum strictly above
every nontrivial separate flattening bound. It demonstrates an exact
finite-word compatibility effect; it is not the long-horizon family
used above. The separation between word-distortion and simplex
enclosure in Theorem~\ref{thm:separation50} likewise distinguishes
two optimizations without substituting either for the full stochastic
problem.

The article proceeds from the machine model and reachable sections to
all-width rigidity, exact simplex geometry and stable occupation. It
then gives the orthogonal finite example, the complete two-state
certificate theory and its encoding bounds. The arithmetic,
word-profile, metric and finite-bit results are retained with their
proofs in the appendices. Section~\ref{sec:resources51} places the
results beside their classical antecedents and states the resource
conventions in one place. The finite-dimensional results do not
import unproved analytic assumptions from other papers in the series.
